
\subsection{算子 $\mathfrak p$、$\mathfrak q$}
\label{subsec:8-10-2}

在定义算子 $\mathfrak p$、$\mathfrak q$ 等时，我们考虑定义域上带内部
标记点的全纯圆盘模空间。

先考察第 $0$ 个标记点是边界标记点的情形。置
\[
 \mathcal M_{(1,k),\ell}(\beta)=\mathcal M_{k+1,\ell}(\beta),
\]
即它是类 $\beta$ 中、带 $k+1$ 个边界标记点和 $\ell$ 个内部标记点的
伪全纯圆盘模空间。用下列等式定义其定向：
\[
\begin{aligned}
 \widetilde{\mathcal M}_{(1,k),\ell}(\beta)
 &:={\widetilde{\mathcal M}}(\beta)
 \times\partial D_0^2\times D_1^2\times\cdots\times D_\ell^2
 \times\partial D_{\ell+1}^2\times\cdots\times\partial D_{\ell+k}^2,\\
 \mathcal M_{(1,k),\ell}(\beta)
 &:={\widetilde{\mathcal M}}_{(1,k),\ell}(\beta)
 /\operatorname{PSL}(2;\mathbb R).
\end{aligned}
\]

\begin{definition}{8.10.1}\label{def:8-10-1}
设 $f_i\colon P_i\to L$ 是 $L$ 中的光滑奇异单形，
$g_j\colon Q_j\to M$ 是 $M$ 中的光滑奇异单形。定义
\[
\begin{aligned}
&\mathcal M_{(1,k),\ell}^{\mathrm{main}}
(\beta;Q_1,\ldots,Q_\ell;P_1,\ldots,P_k)\\
&\quad:={}
(-1)^\epsilon
\mathcal M_{(1,k),\ell}^{\mathrm{main}}(L;\beta)
{}_{\substack{(\operatorname{ev}_1^{\mathrm{int}},\ldots,
\operatorname{ev}_\ell^{\mathrm{int}},\\
\operatorname{ev}_1,\ldots,\operatorname{ev}_k)}}
\mathop{\times}_{\substack{g_1\times\cdots\times g_\ell\\
{}\times f_1\times\cdots\times f_k}}
\left(\prod_{i=1}^{\ell}Q_i\times\prod_{j=1}^{k}P_j\right),
\end{aligned}
\]
其中
\[
 \epsilon=(n+1)\sum_{j=1}^{k-1}\sum_{i=1}^{j}\deg P_i
 +\bigl((k+1)(n+1)+1\bigr)\sum_{i=1}^{\ell}\deg Q_i.
\]
\end{definition}

注意，当 $\ell=0$ 时，模空间 $\mathcal M_{(1,k),0}(\beta)$ 就是
$\mathcal M_{k+1}(\beta)$，且上述定向与定义~\ref{def:8-4-1} 相同。
定义~\ref{def:8-10-1} 只处理模空间的主分支。对其他分支，我们也采用
同一定向约定；这些分支将在下面命题~\ref{prop:8-10-4} 的证明中用到。

当第 $0$ 个标记点是内部标记点时，定义
\[
\begin{aligned}
 \widetilde{\mathcal M}_{k,(1,\ell)}(\beta)
 &:={\widetilde{\mathcal M}}(\beta)
 \times D_0^2\times D_1^2\times\cdots\times D_\ell^2
 \times\partial D_{\ell+1}^2\times\cdots\times\partial D_{\ell+k}^2,\\
 \mathcal M_{k,(1,\ell)}(\beta)
 &:={\widetilde{\mathcal M}}_{k,(1,\ell)}(\beta)
 /\operatorname{PSL}(2;\mathbb R).
\end{aligned}
\]

\begin{definition}{8.10.2}\label{def:8-10-2}
设 $f_i\colon P_i\to L$ 是 $L$ 中的光滑奇异单形，
$g_j\colon Q_j\to M$ 是 $M$ 中的光滑奇异单形。定义
\[
\begin{aligned}
&\mathcal M_{k,(1,\ell)}^{\mathrm{main}}
(\beta;Q_1,\ldots,Q_\ell;P_1,\ldots,P_k)\\
&\quad:={}
(-1)^\epsilon
\mathcal M_{k,(1,\ell)}^{\mathrm{main}}(L;\beta)
{}_{\substack{(\operatorname{ev}_1^{\mathrm{int}},\ldots,
\operatorname{ev}_\ell^{\mathrm{int}},\\
\operatorname{ev}_1,\ldots,\operatorname{ev}_k)}}
\mathop{\times}_{\substack{g_1\times\cdots\times g_\ell\\
{}\times f_1\times\cdots\times f_k}}
\left(\prod_{i=1}^{\ell}Q_i\times\prod_{j=1}^{k}P_j\right),
\end{aligned}
\]
其中
\[
 \epsilon=(n+1)\sum_{j=1}^{k-1}\sum_{i=1}^{j}\deg P_i
 +\bigl((k+1)(n+1)+1\bigr)\sum_{i=1}^{\ell}\deg Q_i.
\]
\end{definition}

若赋予 $\operatorname{PSL}(2;\mathbb R)$ 定向，使嵌入
\[
 g\in\operatorname{PSL}(2;\mathbb R)
 \longmapsto
 \bigl(g\cdot1,g\cdot\sqrt{-1},g\cdot(-1)\bigr)
 \in\partial D_0^2\times\partial D_1^2\times\partial D_2^2
\]
保持定向，则嵌入
\[
 g\in\operatorname{PSL}(2;\mathbb R)
 \longmapsto(g\cdot O,g\cdot z_1)\in D_0^2\times\partial D_1^2
\]
也保持定向。这里 $O$ 是圆盘 $D_0^2$ 的原点，而 $D_0^2$ 由复结构
定向。由于我们对右作用取商，故采用与上述定向相反的定向（见约定
\ref{conv:8-3-1}）。由此，正如式~\eqref{eq:3-8-10-1}，有
\[
 \mathfrak p_{1,0}\equiv i_!\pmod{\Lambda_{0,\mathrm{nov}}^+}.
\]

我们可以把
\[
 \mathcal M_{(1,k),\ell}^{\mathrm{main}}
 (\beta;Q_1,\ldots,Q_\ell;P_1,\ldots,P_k)
\]
与
\[
 \mathcal M_{1,(1,0)}^{\mathrm{main}}
 \bigl(\beta_0;
 \mathcal M_{(1,k),\ell}^{\mathrm{main}}
 (\beta;Q_1,\ldots,Q_\ell;P_1,\ldots,P_k)\bigr)
\]
相识别，其中 $\beta_0$ 是从 $D^2$ 到 $L$ 的常值映射所代表的类。
在这一识别下，不难看出
\[
\begin{aligned}
&\mathcal M_{(1,k),\ell}^{\mathrm{main}}
(\beta;Q_1,\ldots,Q_\ell;P_1,\ldots,P_k)\\
&\quad\subset(-1)^n\partial
\mathcal M_{k,(1,\ell)}^{\mathrm{main}}
(\beta;Q_1,\ldots,Q_\ell;P_1,\ldots,P_k).
\end{aligned}
\]

这里还要指出算子 $\mathfrak p$ 的如下分次对称性。每个内部标记点带有
二维自由度，而 $M$ 的维数为偶数，故对 $i=1,\ldots,\ell-1$
\translatornote{源书在此写“$i\ne j$”，但下式只交换相邻的
$Q_i,Q_{i+1}$，并无指标 $j$；故依公式恢复相邻对换的指标范围。}，显然有
\[
\begin{aligned}
&\mathcal M_{k,(1,\ell)}^{\mathrm{main}}
(\beta;Q_1,\ldots,Q_i,Q_{i+1},\ldots,Q_\ell;P_1,\ldots,P_k)\\
&\quad={}
(-1)^{\deg Q_i\cdot\deg Q_{i+1}}
\mathcal M_{k,(1,\ell)}^{\mathrm{main}}
(\beta;Q_1,\ldots,Q_{i+1},Q_i,\ldots,Q_\ell;P_1,\ldots,P_k),
\end{aligned}
\]
其中
\[
 \deg Q=2n-\dim Q.
\]
不难看出，可以用这些定向约定定义算子 $\mathfrak p$、$\mathfrak q$ 等，
使定理~\ref{thm:3-8-32} 等处关于诸 $Q_i$ 的循环对称性条件带符号成立。

注意
\[
 \dim\mathcal M_{m_1,(1,\ell_1)}^{\mathrm{reg}}(\beta')
 =\dim\mathcal M_{(1,m_1),\ell_1}^{\mathrm{reg}}(\beta')+1
 \equiv n+m_1+1\pmod2.
\]
计及右端的 $+1$，可以如下证明命题~\ref{prop:8-3-3} 的类似结果。
（这个 $+1$ 出现在命题~\ref{prop:8-10-3} 等式右端 $(-1)$ 的指数中。）

\begin{proposition}{8.10.3}\label{prop:8-10-3}
作为带仓西结构的定向空间，有同构
\begingroup
\small
\[
\begin{aligned}
&\partial\mathcal M_{m,(1,\ell)}^{\mathrm{reg}}(\beta'+\beta'')\\
&\quad={}
\bigcup_{\substack{\beta=\beta'+\beta'',\ \,\ell=\ell_1+\ell_2\\
m=m_1+m_2-1}}
(-1)^{(m_1-1)(m_2-1)+(n+m_1)+1}
\begin{gathered}
\mathcal M_{m_1,(1,\ell_1)}^{\mathrm{reg}}(\beta')
{}_{\operatorname{ev}_1^{\beta'}}\\[-2pt]
{}\mathop{\times}_{\operatorname{ev}_0^{\beta''}}
\mathcal M_{(1,m_2),\ell_2}^{\mathrm{reg}}(\beta'')
\end{gathered}.
\end{aligned}
\]
\endgroup
这里并集取遍
$\beta=\beta'+\beta''$、$\ell=\ell_1+\ell_2$ 与
$m=m_1+m_2-1$
\translatornote{源书说明句改用未在本命题出现的 $\beta_1,\beta_2$；
左端、两个因子及评价映射均使用 $\beta',\beta''$，故统一恢复。}。
\end{proposition}

\begin{proposition}{8.10.4}\label{prop:8-10-4}
\begin{enumerate}[label=(\arabic*),leftmargin=2.8em]
\item 对 $\beta=\beta_1+\beta_2$，有
\[
\begin{aligned}
&\mathcal M_{k-k_2+1,(1,\ell_1)}^{\mathrm{main}}
\Bigl(\beta_1;Q_1,\ldots,Q_{\ell_1};P_1,\ldots,P_{i-1},\\
&\qquad
\mathcal M_{(1,k_2),\ell_2}^{\mathrm{main}}
(\beta_2;Q_{\ell_1+1},\ldots,Q_{\ell_1+\ell_2};
P_i,\ldots,P_{i+k_2-1}),P_{i+k_2},\ldots,P_k\Bigr)\\
&\quad\subset(-1)^{\epsilon_1}\partial
\mathcal M_{k,(1,\ell_1+\ell_2)}^{\mathrm{main}}
\bigl(\beta;Q_1,\ldots,Q_{\ell_1},Q_{\ell_1+1},\ldots,Q_{\ell_1+\ell_2};\\
&\hspace{8em}
P_1,\ldots,P_{i-1},P_i,\ldots,P_{i+k_2-1},P_{i+k_2},\ldots,P_k\bigr),
\end{aligned}
\]
其中
\[
\begin{aligned}
 \epsilon_1={}&n+
 \sum_{j=1}^{i-1}(\deg P_j+1)
 +\sum_{j=1}^{\ell_1}\deg Q_j\\
 &+\sum_{j=1}^{i-1}(\deg P_j+1)
 \left(\sum_{h=\ell_1+1}^{\ell_1+\ell_2}\deg Q_h\right).
\end{aligned}
\]

\item 有
\[
\begin{aligned}
&\mathcal M_{k,(1,\ell)}^{\mathrm{main}}
(\beta;Q_1,\ldots,Q_\ell;P_1,\ldots,\partial P_r,\ldots,P_k)\\
&\quad\subset(-1)^{\epsilon_2}\partial
\mathcal M_{k,(1,\ell)}^{\mathrm{main}}
(\beta;Q_1,\ldots,Q_\ell;P_1,\ldots,P_r,\ldots,P_k),
\end{aligned}
\]
其中
\[
 \epsilon_2=\sum_{j=1}^{r-1}(\deg P_j+1)
 +\sum_{j=1}^{\ell}\deg Q_j.
\]

\item 有
\[
\begin{aligned}
&\mathcal M_{k,(1,\ell)}^{\mathrm{main}}
(\beta;Q_1,\ldots,\partial Q_h,\ldots,Q_\ell;P_1,\ldots,P_k)\\
&\quad\subset(-1)^{\epsilon_3}\partial
\mathcal M_{k,(1,\ell)}^{\mathrm{main}}
(\beta;Q_1,\ldots,Q_h,\ldots,Q_\ell;P_1,\ldots,P_k),
\end{aligned}
\]
其中
\[
 \epsilon_3=1+\sum_{j=1}^{h-1}\deg Q_j.
\]
\end{enumerate}
\end{proposition}
\begin{proof}
这是命题~\ref{prop:8-5-1} 证明的一个修改。(1) 的证明分为若干步。

\par\smallskip\noindent\textbf{第 1 步。}
与命题~\ref{prop:8-5-1} 的证明一样，在计算的中间阶段，我们将使用
主分支以外的分支。注意
\[
\begin{aligned}
&\deg\mathcal M_{(1,k_2),\ell_2}^{\mathrm{main}}
(\beta_2;Q_{\ell_1+1},\ldots,Q_{\ell_1+\ell_2};
P_i,\ldots,P_{i+k_2-1})\\
&\quad\equiv
\sum_{j=i}^{i+k_2-1}(\deg P_j+1)
+\sum_{h=\ell_1+1}^{\ell_1+\ell_2}\deg Q_h
\pmod2.
\end{aligned}
\]
反复使用引理~\ref{lem:8-4-3} 的类似结论，得到
\[
\begin{aligned}
&\mathcal M_{k-k_2+1,(1,\ell_1)}^{\mathrm{main}}
\Bigl(\beta_1;Q_1,\ldots,Q_{\ell_1};P_1,\ldots,P_{i-1},\\
&\qquad
\mathcal M_{(1,k_2),\ell_2}^{\mathrm{main}}
(\beta_2;Q_{\ell_1+1},\ldots,Q_{\ell_1+\ell_2};
P_i,\ldots,P_{i+k_2-1}),P_{i+k_2},\ldots,P_k\Bigr)\\
&\quad\subset(-1)^{\sigma_1}
\mathcal M_{k-k_2+1,(1,\ell_1)}
\Bigl(\beta_1;Q_1,\ldots,Q_{\ell_1};\\
&\qquad
\mathcal M_{(1,k_2),\ell_2}^{\mathrm{main}}
(\beta_2;Q_{\ell_1+1},\ldots,Q_{\ell_1+\ell_2};
P_i,\ldots,P_{i+k_2-1}),\\[-2pt]
&\hspace{12em}P_1,\ldots,P_{i-1},P_{i+k_2},\ldots,P_k\Bigr),
\end{aligned}
\]
其中
\[
 \sigma_1=\left(
 1+\sum_{j=i}^{i+k_2-1}(\deg P_j+1)
 +\sum_{h=\ell_1+1}^{\ell_1+\ell_2}\deg Q_h\right)
 \left(\sum_{j=1}^{i-1}(\deg P_j+1)\right).
\]

\par\smallskip\noindent\textbf{第 2 步。}
下面比较
\[
\begin{aligned}
\mathcal M_{k-k_2+1,(1,\ell_1)}\Bigl(&
\beta_1;Q_1,\ldots,Q_{\ell_1};
\mathcal M_{(1,k_2),\ell_2}^{\mathrm{main}}
(\beta_2;Q_{\ell_1+1},\ldots,Q_{\ell_1+\ell_2};\\[-2pt]
&P_i,\ldots,P_{i+k_2-1}),
P_1,\ldots,P_{i-1},P_{i+k_2},\ldots,P_k\Bigr)
\end{aligned}
\]
上的定向与
\[
\begin{aligned}
\partial\mathcal M_{k,(1,\ell_1+\ell_2)}
\bigl(&\beta;Q_1,\ldots,Q_{\ell_1},Q_{\ell_1+1},\ldots,Q_{\ell_1+\ell_2};\\[-2pt]
&P_i,\ldots,P_{i+k_2-1},P_1,\ldots,P_{i-1},
P_{i+k_2},\ldots,P_k\bigr)
\end{aligned}
\]
上的定向。由定义~\ref{def:8-10-2}
\translatornote{源书此处误引定义 8.10.1；本式第 $0$ 个标记点在内部，
其定向由定义 8.10.2 给出。}，有
\begingroup
\small
\[
\begin{aligned}
&\mathcal M_{k,(1,\ell_1+\ell_2)}
(\beta;Q_1,\ldots,Q_{\ell_1+\ell_2};
P_i,\ldots,P_{i+k_2-1},P_1,\ldots,P_{i-1},P_{i+k_2},\ldots,P_k)\\
&\quad={}
(-1)^{\gamma_1}
\mathcal M_{k,(1,\ell_1+\ell_2)}(\beta)_
{(\operatorname{ev}_1^{\mathrm{int}},\ldots,
\operatorname{ev}_{\ell_1+\ell_2}^{\mathrm{int}},
\operatorname{ev}_1,\ldots,\operatorname{ev}_k)}\\[-2pt]
&\hspace{8em}{}\times
\left(\prod_{h=1}^{\ell_1+\ell_2}Q_h
\times\prod_{j=i}^{i+k_2-1}P_j
\times\prod_{j=1}^{i-1}P_j
\times\prod_{j=i+k_2}^{k}P_j\right),
\end{aligned}
\]
\endgroup
其中
\[
\begin{aligned}
\gamma_1={}&(n+1)\Bigg(
\sum_{j=i}^{i+k_2-1}\sum_{\lambda=i}^{j}\deg P_\lambda\\
&+\sum_{j=1}^{i-1}\left(
\sum_{m=i}^{i+k_2-1}\deg P_m
+\sum_{\lambda=1}^{j}\deg P_\lambda\right)\\
&+\sum_{j=i+k_2}^{k-1}\left(
\sum_{m=i}^{i+k_2-1}\deg P_m
+\sum_{m=1}^{i-1}\deg P_m
+\sum_{\lambda=i+k_2}^{j}\deg P_\lambda\right)\Bigg)\\
&+\bigl((k+1)(n+1)+1\bigr)
\sum_{h=1}^{\ell_1+\ell_2}\deg Q_h.
\end{aligned}
\]

使用引理~\ref{lem:8-2-3}(3) 的迭代公式，得到
\begingroup
\small
\[
\begin{aligned}
&\mathcal M_{k,(1,\ell_1+\ell_2)}(\beta)_
{(\operatorname{ev}_1^{\mathrm{int}},\ldots,
\operatorname{ev}_{\ell_1+\ell_2}^{\mathrm{int}},
\operatorname{ev}_1,\ldots,\operatorname{ev}_k)}\\[-2pt]
&\hspace{5em}{}\times
\left(\prod_{h=1}^{\ell_1+\ell_2}Q_h
\times\prod_{j=i}^{i+k_2-1}P_j
\times\prod_{j=1}^{i-1}P_j
\times\prod_{j=i+k_2}^{k}P_j\right)\\
&\quad={}
(-1)^{\gamma_2}
\Bigg(
\mathcal M_{k,(1,\ell_1+\ell_2)}(\beta)_
{(\operatorname{ev}_{\ell_1+1}^{\mathrm{int}},\ldots,
\operatorname{ev}_{\ell_1+\ell_2}^{\mathrm{int}},
\operatorname{ev}_1,\ldots,\operatorname{ev}_{k_2})}\\[-2pt]
&\hspace{9em}{}\times
\left(\prod_{h=\ell_1+1}^{\ell_1+\ell_2}Q_h
\times\prod_{j=i}^{i+k_2-1}P_j\right)
\Bigg)_
{(\operatorname{ev}_1^{\mathrm{int}},\ldots,
\operatorname{ev}_{\ell_1}^{\mathrm{int}},
\operatorname{ev}_{k_2+1},\ldots,\operatorname{ev}_k)}\\[-2pt]
&\hspace{9em}{}\times
\left(\prod_{h=1}^{\ell_1}Q_h
\times\prod_{j=1}^{i-1}P_j
\times\prod_{j=i+k_2}^{k}P_j\right),
\end{aligned}
\]
\endgroup
其中
\[
\begin{aligned}
\gamma_2={}&(k-k_2)n\left(
k_2n+\sum_{j=i}^{i+k_2-1}(n-\deg P_j)
+\sum_{h=\ell_1+1}^{\ell_1+\ell_2}(2n-\deg Q_h)\right)\\
&+\left(\sum_{h=1}^{\ell_1}(2n-\deg Q_h)\right)
\left(\sum_{j=i}^{i+k_2-1}(n-\deg P_j)\right)\\
&+\left(\sum_{h=1}^{\ell_1}(2n-\deg Q_h)\right)
\left(\sum_{h=\ell_1+1}^{\ell_1+\ell_2}(2n-\deg Q_h)\right),
\end{aligned}
\]
并且
\[
\begin{aligned}
\gamma_2\equiv{}&
(k-k_2)n\sum_{j=i}^{i+k_2-1}\deg P_j
+(k-k_2)n\sum_{h=\ell_1+1}^{\ell_1+\ell_2}\deg Q_h\\
&+\left(\sum_{h=1}^{\ell_1}\deg Q_h\right)
 \left(\sum_{j=i}^{i+k_2-1}(n-\deg P_j)\right)\\
&+\left(\sum_{h=1}^{\ell_1}\deg Q_h\right)
 \left(\sum_{h=\ell_1+1}^{\ell_1+\ell_2}\deg Q_h\right)
\pmod2.
\end{aligned}
\]

由命题~\ref{prop:8-10-3}，得到
\begingroup
\small
\[
\begin{aligned}
&\partial\Bigg[
\Bigg(
\mathcal M_{k,(1,\ell_1+\ell_2)}(\beta)_
{(\operatorname{ev}_{\ell_1+1}^{\mathrm{int}},\ldots,
\operatorname{ev}_{\ell_1+\ell_2}^{\mathrm{int}},
\operatorname{ev}_1,\ldots,\operatorname{ev}_{k_2})}\\[-2pt]
&\hspace{9em}{}\times
\left(\prod_{h=\ell_1+1}^{\ell_1+\ell_2}Q_h
\times\prod_{j=i}^{i+k_2-1}P_j\right)
\Bigg)_
{(\operatorname{ev}_1^{\mathrm{int}},\ldots,
\operatorname{ev}_{\ell_1}^{\mathrm{int}},
\operatorname{ev}_{k_2+1},\ldots,\operatorname{ev}_k)}\\[-2pt]
&\hspace{9em}{}\times
\left(\prod_{h=1}^{\ell_1}Q_h
\times\prod_{j=1}^{i-1}P_j
\times\prod_{j=i+k_2}^{k}P_j\right)\Bigg]\\
&\quad\supset(-1)^{\gamma_3}
\Bigg[
\Bigg(
\mathcal M_{k-k_2+1,(1,\ell_1)}(\beta_1)_
{\operatorname{ev}_1^{\beta_1}}
\mathop{\times}_{\operatorname{ev}_0^{\beta_2}}
\mathcal M_{(1,k_2),\ell_2}(\beta_2)
\Bigg)_
{(\operatorname{ev}_{\ell_1+1}^{\beta,\mathrm{int}},\ldots,
\operatorname{ev}_{\ell_1+\ell_2}^{\beta,\mathrm{int}},
\operatorname{ev}_1^\beta,\ldots,\operatorname{ev}_{k_2}^\beta)}\\[-2pt]
&\hspace{9em}{}\times
\left(\prod_{h=\ell_1+1}^{\ell_1+\ell_2}Q_h
\times\prod_{j=i}^{i+k_2-1}P_j\right)
\Bigg]_
{(\operatorname{ev}_1^{\beta,\mathrm{int}},\ldots,
\operatorname{ev}_{\ell_1}^{\beta,\mathrm{int}},
\operatorname{ev}_{k_2+1}^\beta,\ldots,\operatorname{ev}_k^\beta)}\\[-2pt]
&\hspace{9em}{}\times
\left(\prod_{h=1}^{\ell_1}Q_h
\times\prod_{j=1}^{i-1}P_j
\times\prod_{j=i+k_2}^{k}P_j\right),
\end{aligned}
\]
\endgroup
其中
\[
 \gamma_3=(k-k_2)(k_2-1)+n+k-k_2
 \equiv n+kk_2+k_2\pmod2.
\]

由引理~\ref{lem:8-2-3}(2) 的结合性质与定义~\ref{def:8-10-1}，有
\begingroup
\small
\[
\begin{aligned}
&\Bigg[
\Bigg(
\mathcal M_{k-k_2+1,(1,\ell_1)}(\beta_1)_
{\operatorname{ev}_1^{\beta_1}}
\mathop{\times}_{\operatorname{ev}_0^{\beta_2}}
\mathcal M_{(1,k_2),\ell_2}(\beta_2)
\Bigg)_
{(\operatorname{ev}_{\ell_1+1}^{\beta,\mathrm{int}},\ldots,
\operatorname{ev}_{\ell_1+\ell_2}^{\beta,\mathrm{int}},
\operatorname{ev}_1^\beta,\ldots,\operatorname{ev}_{k_2}^\beta)}\\[-2pt]
&\hspace{8em}{}\times
\left(\prod_{h=\ell_1+1}^{\ell_1+\ell_2}Q_h
\times\prod_{j=i}^{i+k_2-1}P_j\right)
\Bigg]_
{(\operatorname{ev}_1^{\beta,\mathrm{int}},\ldots,
\operatorname{ev}_{\ell_1}^{\beta,\mathrm{int}},
\operatorname{ev}_{k_2+1}^\beta,\ldots,\operatorname{ev}_k^\beta)}\\[-2pt]
&\hspace{8em}{}\times
\left(\prod_{h=1}^{\ell_1}Q_h
\times\prod_{j=1}^{i-1}P_j
\times\prod_{j=i+k_2}^{k}P_j\right)\\
&={}
\Bigg[
\mathcal M_{k-k_2+1,(1,\ell_1)}(\beta_1)_
{\operatorname{ev}_1^{\beta_1}}
\mathop{\times}_{\operatorname{ev}_0^{\beta_2}}
\Bigg(
\mathcal M_{(1,k_2),\ell_2}(\beta_2)_
{(\operatorname{ev}_1^{\beta_2,\mathrm{int}},\ldots,
\operatorname{ev}_{\ell_2}^{\beta_2,\mathrm{int}},
\operatorname{ev}_1^{\beta_2},\ldots,
\operatorname{ev}_{k_2}^{\beta_2})}\\[-2pt]
&\hspace{13em}{}\times
\left(\prod_{h=\ell_1+1}^{\ell_1+\ell_2}Q_h
\times\prod_{j=i}^{i+k_2-1}P_j\right)
\Bigg)\Bigg]_
{\substack{(\operatorname{ev}_1^{\beta,\mathrm{int}},\ldots,\\
\operatorname{ev}_{\ell_1}^{\beta,\mathrm{int}},
\operatorname{ev}_{k_2+1}^\beta,\ldots,\\
\operatorname{ev}_k^\beta)}}\\[-2pt]
&\hspace{8em}{}\times
\left(\prod_{h=1}^{\ell_1}Q_h
\times\prod_{j=1}^{i-1}P_j
\times\prod_{j=i+k_2}^{k}P_j\right)\\
&={}(-1)^{\gamma_4}
\left(
\begin{gathered}
\mathcal M_{k-k_2+1,(1,\ell_1)}(\beta_1)_
{\operatorname{ev}_1^{\beta_1}}\\[-2pt]
{}\mathop{\times}_{\operatorname{ev}_0^{\beta_2}}
\mathcal M_{(1,k_2),\ell_2}\\[-2pt]
(\beta_2;Q_{\ell_1+1},\ldots,Q_{\ell_1+\ell_2};\\[-2pt]
P_i,\ldots,P_{i+k_2-1})
\end{gathered}
\right)_
{\substack{(\operatorname{ev}_1^{\beta,\mathrm{int}},\ldots,\\
\operatorname{ev}_{\ell_1}^{\beta,\mathrm{int}},
\operatorname{ev}_{k_2+1}^\beta,\ldots,\\
\operatorname{ev}_k^\beta)}}\\[-2pt]
&\hspace{8em}{}\times
\left(\prod_{h=1}^{\ell_1}Q_h
\times\prod_{j=1}^{i-1}P_j
\times\prod_{j=i+k_2}^{k}P_j\right),
\end{aligned}
\]
\endgroup
其中
\[
 \gamma_4=(n+1)
 \sum_{j=i}^{i+k_2-2}\sum_{\lambda=i}^{j}\deg P_\lambda
 +\bigl((k_2+1)(n+1)+1\bigr)
 \sum_{h=\ell_1+1}^{\ell_1+\ell_2}\deg Q_h.
\]

再次使用迭代公式和定义~\ref{def:8-10-2}，得到
\begingroup
\small
\[
\begin{aligned}
&\left(
\begin{gathered}
\mathcal M_{k-k_2+1,(1,\ell_1)}(\beta_1)_
{\operatorname{ev}_1^{\beta_1}}\\[-2pt]
{}\mathop{\times}_{\operatorname{ev}_0^{\beta_2}}
\mathcal M_{(1,k_2),\ell_2}\\[-2pt]
(\beta_2;Q_{\ell_1+1},\ldots,Q_{\ell_1+\ell_2};
P_i,\ldots,P_{i+k_2-1})
\end{gathered}
\right)_
{\substack{(\operatorname{ev}_1^{\beta,\mathrm{int}},\ldots,
\operatorname{ev}_{\ell_1}^{\beta,\mathrm{int}},\\
\operatorname{ev}_{k_2+1}^\beta,\ldots,\operatorname{ev}_k^\beta)}}\\[-2pt]
&\hspace{8em}{}\times
\left(\prod_{h=1}^{\ell_1}Q_h
\times\prod_{j=1}^{i-1}P_j
\times\prod_{j=i+k_2}^{k}P_j\right)\\
&={}
(-1)^{\gamma_5}
\mathcal M_{k-k_2+1,(1,\ell_1)}(\beta_1)_
{(\operatorname{ev}_1^{\beta_1,\mathrm{int}},\ldots,
\operatorname{ev}_{\ell_1}^{\beta_1,\mathrm{int}},
\operatorname{ev}_1^{\beta_1},\ldots,
\operatorname{ev}_{k-k_2+1}^{\beta_1})}\\[-2pt]
&\hspace{5em}{}\times\Bigg(
\mathcal M_{(1,k_2),\ell_2}
(\beta_2;Q_{\ell_1+1},\ldots,Q_{\ell_1+\ell_2};
P_i,\ldots,P_{i+k_2-1})\\[-2pt]
&\hspace{14em}{}\times
\left(\prod_{h=1}^{\ell_1}Q_h
\times\prod_{j=1}^{i-1}P_j
\times\prod_{j=i+k_2}^{k}P_j\right)\Bigg)\\
&={}
(-1)^{\gamma_5+\gamma_6}
\mathcal M_{k-k_2+1,(1,\ell_1)}
\Bigl(\beta_1;Q_1,\ldots,Q_{\ell_1};\\[-2pt]
&\hspace{6em}{}
\mathcal M_{(1,k_2),\ell_2}
(\beta_2;Q_{\ell_1+1},\ldots,Q_{\ell_1+\ell_2};
P_i,\ldots,P_{i+k_2-1}),\\[-2pt]
&\hspace{6em}{}
P_1,\ldots,P_{i-1},P_{i+k_2},\ldots,P_k\Bigr),
\end{aligned}
\]
\endgroup
其中
\[
\begin{aligned}
\gamma_5
&=n(k-k_2)\left(
2n-\sum_{j=i}^{i+k_2-1}(\deg P_j+1)
-\sum_{h=\ell_1+1}^{\ell_1+\ell_2}\deg Q_h\right)\\
&\equiv n(k-k_2)\left(
k_2+\sum_{j=i}^{i+k_2-1}\deg P_j
-\sum_{h=\ell_1+1}^{\ell_1+\ell_2}\deg Q_h\right)
\pmod2,
\end{aligned}
\]
而
\begingroup
\small
\[
\begin{aligned}
\gamma_6={}&(n+1)\Bigg(
\deg\mathcal M_{(1,k_2),\ell_2}
(\beta_2;Q_{\ell_1+1},\ldots,Q_{\ell_1+\ell_2};
P_i,\ldots,P_{i+k_2-1})\\
&+\sum_{j=1}^{i-1}\Bigg(
\deg\mathcal M_{(1,k_2),\ell_2}
(\beta_2;Q_{\ell_1+1},\ldots,Q_{\ell_1+\ell_2};
P_i,\ldots,P_{i+k_2-1})
+\sum_{\lambda=1}^{j}\deg P_\lambda\Bigg)\\
&+\sum_{j=i+k_2}^{k-1}\Bigg(
\deg\mathcal M_{(1,k_2),\ell_2}
(\beta_2;Q_{\ell_1+1},\ldots,Q_{\ell_1+\ell_2};
P_i,\ldots,P_{i+k_2-1})\\[-2pt]
&\hspace{12em}{}+\sum_{m=1}^{i-1}\deg P_m
+\sum_{\lambda=i+k_2}^{j}\deg P_\lambda\Bigg)\Bigg)\\
&+\bigl((k-k_2+2)(n+1)+1\bigr)
\sum_{h=1}^{\ell_1}\deg Q_h\\
&+\dim\mathcal M_{(1,k_2),\ell_2}
(\beta_2;Q_{\ell_1+1},\ldots,Q_{\ell_1+\ell_2};
P_i,\ldots,P_{i+k_2-1})
\sum_{h=1}^{\ell_1}(2n-\deg Q_h).
\end{aligned}
\]
\endgroup
这里已把 $\gamma_6$ 前两项的内部输入起点恢复为 $Q_{\ell_1+1}$，
且上一串等式最后一行的外层类保持为 $\beta_1$
\translatornote{源书在 $\gamma_6$ 的前两项把输入起点误写成
$Q_{\ell_1}$，却在后一项及本证明其余各处写 $Q_{\ell_1+1}$；又把上一串
等式最后一行的外层类由 $\beta_1$ 误写成总类 $\beta$。这里依命题的
$\beta=\beta_1+\beta_2$ 分解与输入分块统一恢复。}。

模 $2$，$\gamma_6$ 同余于
\begingroup
\small
\[
\begin{aligned}
\gamma_6\equiv{}&(n+1)\Bigg(
(k-k_2)\left(k_2+
\sum_{m=i}^{i+k_2-1}\deg P_m
+\sum_{h=\ell_1+1}^{\ell_1+\ell_2}\deg Q_h\right)\\
&+(k-i-k_2)\left(\sum_{m=1}^{i-1}\deg P_m\right)
+\sum_{j=1}^{i-1}\sum_{\lambda=1}^{j}\deg P_\lambda
+\sum_{j=i+k_2}^{k-1}\sum_{\lambda=i+k_2}^{j}\deg P_\lambda
\Bigg)\\
&+\bigl((k-k_2)(n+1)+1\bigr)
\sum_{h=1}^{\ell_1}\deg Q_h\\
&+\left(\sum_{h=1}^{\ell_1}\deg Q_h\right)
\left(n-\left(k_2+
\sum_{m=i}^{i+k_2-1}\deg P_m
+\sum_{h=\ell_1+1}^{\ell_1+\ell_2}\deg Q_h\right)\right)
\pmod2.
\end{aligned}
\]
\endgroup
于是，初等计算给出
\[
 \gamma_1+\gamma_2+\gamma_3+\gamma_4+\gamma_5+\gamma_6
 \equiv n+\sum_{h=1}^{\ell_1}\deg Q_h\pmod2.
\]
因而
\[
\begin{aligned}
&\mathcal M_{k-k_2+1,(1,\ell_1)}
\Bigl(\beta_1;Q_1,\ldots,Q_{\ell_1};\\[-2pt]
&\hspace{4em}{}
\mathcal M_{(1,k_2),\ell_2}^{\mathrm{main}}
(\beta_2;Q_{\ell_1+1},\ldots,Q_{\ell_1+\ell_2};
P_i,\ldots,P_{i+k_2-1}),\\[-2pt]
&\hspace{4em}{}
P_1,\ldots,P_{i-1},P_{i+k_2},\ldots,P_k\Bigr)\\
&\quad\subset
(-1)^{n+\sum_{h=1}^{\ell_1}\deg Q_h}
\partial\mathcal M_{k,(1,\ell_1+\ell_2)}
\bigl(\beta;Q_1,\ldots,Q_{\ell_1+\ell_2};\\[-2pt]
&\hspace{5em}{}
P_i,\ldots,P_{i+k_2-1},P_1,\ldots,P_{i-1},
P_{i+k_2},\ldots,P_k\bigr).
\end{aligned}
\]

\par\smallskip\noindent\textbf{第 3 步。}
另一方面，再次使用引理~\ref{lem:8-4-3} 可见
\[
\begin{aligned}
&\mathcal M_{k,(1,\ell_1+\ell_2)}
(\beta;Q_1,\ldots,Q_{\ell_1+\ell_2};\\[-2pt]
&\hspace{4em}{}
P_i,\ldots,P_{i+k_2-1},P_1,\ldots,P_{i-1},
P_{i+k_2},\ldots,P_k)\\
&\quad={}
(-1)^{\sigma_2}
\mathcal M_{k,(1,\ell_1+\ell_2)}^{\mathrm{main}}
(\beta;Q_1,\ldots,Q_{\ell_1+\ell_2};\\[-2pt]
&\hspace{4em}{}
P_1,\ldots,P_{i-1},P_i,\ldots,P_{i+k_2-1},
P_{i+k_2},\ldots,P_k),
\end{aligned}
\]
其中
\[
 \sigma_2=\left(\sum_{j=i}^{i+k_2-1}(\deg P_j+1)\right)
 \left(\sum_{j=1}^{i-1}(\deg P_j+1)\right).
\]
因此
\[
\begin{aligned}
&\sigma_1+n+\sum_{h=1}^{\ell_1}\deg Q_h+\sigma_2\\
&\quad\equiv n+\sum_{h=1}^{\ell_1}\deg Q_h
+\sum_{j=1}^{i-1}(\deg P_j+1)\\
&\qquad{}+
\sum_{j=1}^{i-1}(\deg P_j+1)
\left(\sum_{h=\ell_1+1}^{\ell_1+\ell_2}\deg Q_h\right)
\pmod2,
\end{aligned}
\]
这就证明了命题~\ref{prop:8-10-4}(1)。

下面证明 (2) 与 (3)。回忆
\[
\begin{aligned}
&\mathcal M_{k,(1,\ell)}^{\mathrm{main}}
(\beta;Q_1,\ldots,Q_\ell;P_1,\ldots,P_i,\ldots,P_k)\\
&\quad={}
(-1)^{\delta_1}
\mathcal M_{k,(1,\ell)}^{\mathrm{main}}(\beta)_
{\substack{(\operatorname{ev}_1^{\mathrm{int}},\ldots,
\operatorname{ev}_\ell^{\mathrm{int}},\\
\operatorname{ev}_1,\ldots,\operatorname{ev}_k)}}
\mathop{\times}_{\substack{g_1\times\cdots\times g_\ell\\
{}\times f_1\times\cdots\times f_k}}
\left(\prod_{h=1}^{\ell}Q_h\times\prod_{j=1}^{k}P_j\right),
\end{aligned}
\]
其中
\[
 \delta_1=(n+1)\sum_{j=1}^{k-1}\sum_{i=1}^{j}\deg P_i
 +\bigl((k+1)(n+1)+1\bigr)\sum_{i=1}^{\ell}\deg Q_i.
\]
由引理~\ref{lem:8-2-3}(1)，有
\begingroup
\small
\[
\begin{aligned}
&\partial\Bigg(
\mathcal M_{k,(1,\ell)}^{\mathrm{main}}(\beta)_
{\substack{(\operatorname{ev}_1^{\mathrm{int}},\ldots,
\operatorname{ev}_\ell^{\mathrm{int}},\\
\operatorname{ev}_1,\ldots,\operatorname{ev}_k)}}
\mathop{\times}_{\substack{g_1\times\cdots\times g_\ell\\
{}\times f_1\times\cdots\times f_k}}
\left(\prod_{h=1}^{\ell}Q_h\times\prod_{j=1}^{k}P_j\right)\Bigg)\\
&={}
\partial\mathcal M_{k,(1,\ell)}^{\mathrm{main}}(\beta)_
{\substack{(\operatorname{ev}_1^{\mathrm{int}},\ldots,
\operatorname{ev}_\ell^{\mathrm{int}},\\
\operatorname{ev}_1,\ldots,\operatorname{ev}_k)}}
\mathop{\times}_{\substack{g_1\times\cdots\times g_\ell\\
{}\times f_1\times\cdots\times f_k}}
\left(\prod_{h=1}^{\ell}Q_h\times\prod_{j=1}^{k}P_j\right)\\
&\quad\sqcup(-1)^{n+k+1+nk}
\mathcal M_{k,(1,\ell)}^{\mathrm{main}}(\beta)_
{\substack{(\operatorname{ev}_1^{\mathrm{int}},\ldots,
\operatorname{ev}_\ell^{\mathrm{int}},\\
\operatorname{ev}_1,\ldots,\operatorname{ev}_k)}}
\mathop{\times}_{\substack{g_1\times\cdots\times g_\ell\\
{}\times f_1\times\cdots\times f_k}}
\partial\left(\prod_{h=1}^{\ell}Q_h\times\prod_{j=1}^{k}P_j\right),
\end{aligned}
\]
\endgroup
因为
\[
 \dim\mathcal M_{k,(1,\ell)}^{\mathrm{main}}(\beta)
 =n+\mu(\beta)-3+k+2(\ell+1)
 \equiv n+k+1\pmod2,
\]
且 $M$ 为偶数维。还要注意
\begingroup
\small
\[
\begin{aligned}
&\partial\left(\prod_{i=1}^{\ell}Q_i\times\prod_{j=1}^{k}P_j\right)\\
&={}
\coprod_{h=1}^{\ell}
(-1)^{\sum_{i=1}^{h-1}\dim Q_i}
Q_1\times\cdots\times\partial Q_h\times\cdots\times Q_\ell
\times\prod_{j=1}^{k}P_j\\
&\quad\sqcup
\coprod_{r=1}^{k}
(-1)^{\sum_{i=1}^{\ell}\dim Q_i+
\sum_{j=1}^{r-1}\dim P_j}
\prod_{i=1}^{\ell}Q_i\times
P_1\times\cdots\times\partial P_r\times\cdots\times P_k.
\end{aligned}
\]
\endgroup

另一方面，定义~\ref{def:8-10-2} 给出
\begingroup
\small
\[
\begin{aligned}
&\mathcal M_{k,(1,\ell)}^{\mathrm{main}}
(\beta;Q_1,\ldots,Q_\ell;P_1,\ldots,\partial P_r,\ldots,P_k)\\
&\quad={}
(-1)^{\delta_2}
\mathcal M_{k,(1,\ell)}^{\mathrm{main}}(\beta)_
{(\operatorname{ev}_1^{\mathrm{int}},\ldots,
\operatorname{ev}_\ell^{\mathrm{int}},
\operatorname{ev}_1,\ldots,\operatorname{ev}_k)}
\mathop{\times}_{g_1\times\cdots\times g_\ell
\times f_1\times\cdots\times f_k}\\[-2pt]
&\hspace{11em}{}
\left(\prod_{i=1}^{\ell}Q_i\times
\prod_{j=1}^{r-1}P_j\times\partial P_r\times
\prod_{j=r+1}^{k}P_j\right),
\end{aligned}
\]
\endgroup
其中
\[
 \delta_2=(n+1)
 \sum_{j=1}^{k-1}\sum_{\lambda=1}^{j}(\deg P_\lambda)'
 +\bigl((k+1)(n+1)+1\bigr)\sum_{j=1}^{\ell}\deg Q_j,
\]
这里
\[
 (\deg P_\lambda)'=
 \begin{cases}
 \deg P_\lambda,&\lambda\ne r,\\
 \deg P_\lambda+1,&\lambda=r.
 \end{cases}
\]
因此
\[
\begin{aligned}
&\delta_1+n+k+1+nk
+\sum_{j=1}^{r-1}(n-\deg P_j)
+\sum_{i=1}^{\ell}(2n-\deg Q_i)+\delta_2\\
&\quad\equiv
\sum_{j=1}^{r-1}(\deg P_j+1)
+\sum_{i=1}^{\ell}\deg Q_i\pmod2,
\end{aligned}
\]
这就证明了命题~\ref{prop:8-10-4}(2)。

再回忆
\begingroup
\small
\[
\begin{aligned}
&\mathcal M_{k,(1,\ell)}^{\mathrm{main}}
(\beta;Q_1,\ldots,\partial Q_h,\ldots,Q_\ell;P_1,\ldots,P_k)\\
&\quad={}
(-1)^{\delta_3}
\mathcal M_{k,(1,\ell)}^{\mathrm{main}}(\beta)_
{(\operatorname{ev}_1^{\mathrm{int}},\ldots,
\operatorname{ev}_\ell^{\mathrm{int}},
\operatorname{ev}_1,\ldots,\operatorname{ev}_k)}
\mathop{\times}_{g_1\times\cdots\times g_\ell
\times f_1\times\cdots\times f_k}\\[-2pt]
&\hspace{9em}{}
\left(\prod_{i=1}^{h-1}Q_i\times\partial Q_h\times
\prod_{i=h+1}^{\ell}Q_i\times\prod_{j=1}^{k}P_j\right),
\end{aligned}
\]
\endgroup
其中
\[
 \delta_3=(n+1)\sum_{j=1}^{k-1}\sum_{i=1}^{j}\deg P_i
 +\bigl((k+1)(n+1)+1\bigr)
 \left(\sum_{i=1}^{\ell}\deg Q_i-1\right).
\]
因此
\[
 \delta_1+n+k+1+nk
 +\sum_{i=1}^{h-1}(2n-\deg Q_i)+\delta_3
 \equiv\sum_{i=1}^{h-1}\deg Q_i+1\pmod2,
\]
这就证明了命题~\ref{prop:8-10-4}(3)。
\end{proof}

用同样的方法得到下述结果。

\begin{proposition}{8.10.5}\label{prop:8-10-5}
\begin{enumerate}[label=(\arabic*),leftmargin=2.8em]
\item 对 $\beta=\beta_1+\beta_2$，有
\[
\begin{aligned}
&\mathcal M_{(1,k-k_2+1),\ell_1}^{\mathrm{main}}
\Bigl(\beta_1;Q_1,\ldots,Q_{\ell_1};P_1,\ldots,P_{i-1},\\
&\qquad
\mathcal M_{(1,k_2),\ell_2}^{\mathrm{main}}
(\beta_2;Q_{\ell_1+1},\ldots,Q_{\ell_1+\ell_2};
P_i,\ldots,P_{i+k_2-1}),P_{i+k_2},\ldots,P_k\Bigr)\\
&\quad\subset(-1)^{\epsilon'_1}\partial
\mathcal M_{(1,k),\ell_1+\ell_2}^{\mathrm{main}}
\bigl(\beta;Q_1,\ldots,Q_{\ell_1},Q_{\ell_1+1},\ldots,Q_{\ell_1+\ell_2};\\
&\hspace{8em}
P_1,\ldots,P_{i-1},P_i,\ldots,P_{i+k_2-1},P_{i+k_2},\ldots,P_k\bigr),
\end{aligned}
\]
其中
\[
\begin{aligned}
 \epsilon'_1={}&n+1+
 \sum_{j=1}^{i-1}(\deg P_j+1)
 +\sum_{j=1}^{\ell_1}\deg Q_j\\
 &+\sum_{j=1}^{i-1}(\deg P_j+1)
 \left(\sum_{h=\ell_1+1}^{\ell_1+\ell_2}\deg Q_h\right).
\end{aligned}
\]

\item 有
\[
\begin{aligned}
&\mathcal M_{(1,k),\ell}^{\mathrm{main}}
(\beta;Q_1,\ldots,Q_\ell;P_1,\ldots,\partial P_r,\ldots,P_k)\\
&\quad\subset(-1)^{\epsilon'_2}\partial
\mathcal M_{(1,k),\ell}^{\mathrm{main}}
(\beta;Q_1,\ldots,Q_\ell;P_1,\ldots,P_r,\ldots,P_k),
\end{aligned}
\]
其中
\[
 \epsilon'_2=\sum_{j=1}^{r-1}(\deg P_j+1)
 +\sum_{j=1}^{\ell}\deg Q_j+1.
\]

\item 有
\[
\begin{aligned}
&\mathcal M_{(1,k),\ell}^{\mathrm{main}}
(\beta;Q_1,\ldots,\partial Q_h,\ldots,Q_\ell;P_1,\ldots,P_k)\\
&\quad\subset(-1)^{\epsilon'_3}\partial
\mathcal M_{(1,k),\ell}^{\mathrm{main}}
(\beta;Q_1,\ldots,Q_h,\ldots,Q_\ell;P_1,\ldots,P_k),
\end{aligned}
\]
其中
\[
 \epsilon'_3=\sum_{j=1}^{h-1}\deg Q_j.
\]
\end{enumerate}
\end{proposition}

下面核对式~\eqref{eq:3-8-10-3} 中的符号。回忆，由式
\eqref{eq:3-8-1}，
\[
 GW_{0,1}(M;\widetilde\beta)(L)
 =\operatorname{ev}_{0*}\left(
 \mathcal M_{0,2}(M;\widetilde\beta)_{
 \operatorname{ev}_1}\mathbin{\times_M}L\right),
\]
其中 $\widetilde\beta\in H_2(M)$，而 $\beta$ 是它在映射
$i_*\colon H_2(M)\to H_2(M,L)$ 下的像；见第~\ref{subsec:3-8-1} 小节。
我们需要考察 $\mathcal M_{0,2}(M;\widetilde\beta)$ 与
$\mathcal M_{1,(1,0)}(L;\beta)$ 的定向。下述命题足以核对
式~\eqref{eq:3-8-10-3} 中的符号。

\begin{proposition}{8.10.6}\label{prop:8-10-6}
纤维积
\[
 \mathcal M_{0,2}^{\mathrm{reg}}(M;\widetilde\beta)_{
 \operatorname{ev}_1}\mathbin{\times_M}L
\]
的定向，与把它看作
$(-1)^n\mathcal M_{1,(1,0)}^{\mathrm{reg}}(L;\beta)$ 的边界所得的定向
一致
\translatornote{源书在命题陈述及证明结论中把后一个模空间的首下标
写成 $0$；但证明固定了第一个边界标记点，并在全部切空间识别中使用
$\mathcal M_{1,(1,0)}^{\mathrm{reg}}$。依定义 8.10.2，只有首下标为
$1$ 时该边界评价点才存在，故两处均恢复为 $1$。}。
\end{proposition}

\begin{proof}
固定 $0,\infty\in\mathbb{CP}^1$ 为第 $0$、第 $1$ 个标记点，并写成
\[
 \widetilde{\mathcal M}^{\mathrm{reg}}(M;\widetilde\beta)
 /\operatorname{Aut}(\mathbb{CP}^1;0,\infty)
 =\mathcal M_{0,2}^{\mathrm{reg}}(M;\widetilde\beta).
\]
类似地，在闭单位圆盘 $D^2$ 上固定 $0,1$，分别作为第 $0$ 个内部
标记点和第 $1$ 个边界标记点。于是，模空间
$\mathcal M_{1,(1,0)}^{\mathrm{reg}}(L;0)$（相应地，
$\mathcal M_{1,(1,0)}^{\mathrm{reg}}(L;\beta)$）与从 $D^2$ 出发、
代表类 $0$（相应地，$\beta$）的伪全纯映射空间
$\widetilde{\mathcal M}^{\mathrm{reg}}(L;0)$（相应地，
$\widetilde{\mathcal M}^{\mathrm{reg}}(L;\beta)$）相识别。

与引理~\ref{lem:8-3-5} 类似，可知胶合映射
\[
 \widetilde{\mathcal M}^{\mathrm{reg}}(M;\widetilde\beta)_{
 \operatorname{ev}_1}
 \mathop{\times}_{\operatorname{ev}_0^{\mathrm{int}}}
 \widetilde{\mathcal M}^{\mathrm{reg}}(L;0)
 \longrightarrow
 \widetilde{\mathcal M}^{\mathrm{reg}}(L;\beta)
\]
在仓西结构意义下保持定向。因此，在切空间层次有识别
\[
 \left(
 \mathcal M_{0,2}^{\mathrm{reg}}(M;\widetilde\beta)
 \times\operatorname{Aut}(\mathbb{CP}^1;0,\infty)
 \right)_{\operatorname{ev}_1}
 \mathop{\times}_{\operatorname{ev}_0^{\mathrm{int}}}
 \mathcal M_{1,(1,0)}^{\mathrm{reg}}(L;0)
 =\mathcal M_{1,(1,0)}^{\mathrm{reg}}(L;\beta).
\]
注意
\[
 \operatorname{Aut}(\mathbb{CP}^1;0,\infty)
 \cong\mathbb C^\times\cong\mathbb R_{>0}\times S^1.
\]
赋予 $\mathbb R_{>0}$ 定向，使它在 $\mathbb{CP}^1$ 上的作用诱导从
$0$ 到 $\infty$ 的流。于是，与引理~\ref{lem:8-3-5} 的证明一样，
$\mathbb R_{>0}$ 对应于外法向量场。还要注意，
\[
 \operatorname{ev}_0^{\mathrm{int}}\colon
 \mathcal M_{1,(1,0)}^{\mathrm{reg}}(L;0)\longrightarrow L
\]
保持定向。因此
\[
 \mathbb R_{\mathrm{out}}\times
 \left(\mathcal M_{0,2}^{\mathrm{reg}}(M;\widetilde\beta)
 \times S^1\right)_{\operatorname{ev}_1}
 \mathbin{\times_M}L
 =\mathcal M_{1,(1,0)}^{\mathrm{reg}}(L;\beta).
\]
注意，按照式~\eqref{eq:8-2-1-2}，必须把上式左端的 $S^1$ 移到最后一个
因子
\translatornote{源书误引不存在的“约定 (8.2.2.1)”；右作用取商时群因子
置于末位的规则是式 (8.2.1.2)，故据此改引。}。再由式
\eqref{eq:8-2-1-1}，得到
\[
 \mathcal M_{0,2}^{\mathrm{reg}}(M;\widetilde\beta)_{
 \operatorname{ev}_1}\mathbin{\times_M}L
 \subset(-1)^n\partial
 \mathcal M_{1,(1,0)}^{\mathrm{reg}}(L;\beta).
\]
\end{proof}

最后证明引理~\ref{lem:3-8-25} 中公式
\[
 \mathfrak p_{k,\beta}(P_1,\ldots,P_k)
 =(-1)^{\deg'P_k\cdot
 (\sum_{i=1}^{k-1}\deg'P_i)}
 \mathfrak p_{k,\beta}(P_k,P_1,\ldots,P_{k-1})
\]
的符号。这里 $\deg'$ 是移位次数。

下述结果是引理~\ref{lem:8-4-3} 的直接推广。

\begin{lemma}{8.10.7}\label{lem:8-10-7}
设 $\sigma$ 是第 $k$ 个对称群 $\mathfrak S_k$ 中的对换
$(i,i+1)$（$i=1,\ldots,k-1$）。通过改变标记点次序，$\sigma$ 在
\[
 \mathcal M_{k,(1,\ell)}
 (\beta;Q_1,\ldots,Q_\ell;P_1,\ldots,P_i,P_{i+1},\ldots,P_k)
\]
上的作用由下式描述：
\[
\begin{aligned}
&\sigma\Bigl(
\mathcal M_{k,(1,\ell)}
(\beta;Q_1,\ldots,Q_\ell;P_1,\ldots,P_i,P_{i+1},\ldots,P_k)
\Bigr)\\
&\quad={}
(-1)^{(\deg P_i+1)(\deg P_{i+1}+1)}
\mathcal M_{k,(1,\ell)}
(\beta;Q_1,\ldots,Q_\ell;P_1,\ldots,P_{i+1},P_i,\ldots,P_k).
\end{aligned}
\]
\end{lemma}

\begin{proof}
回忆定义~\ref{def:8-10-2}。于是可用与引理~\ref{lem:8-4-3}
相同的方法证明本引理。
\end{proof}

令 $\ell=0$，引理~\ref{lem:8-10-7} 便立即推出引理
\ref{lem:3-8-25} 中所需的符号。
